
\documentclass{article}
\usepackage{colortbl}
\usepackage{makecell}
\usepackage{multirow}
\usepackage{supertabular}

\begin{document}

\newcounter{utterance}

\centering \large Interaction Transcript for game `codenames', experiment `opponent\_difficult', episode 0 with gemma{-}4{-}26b{-}a4b{-}it.
\vspace{24pt}

{ \footnotesize  \setcounter{utterance}{1}
\setlength{\tabcolsep}{0pt}
\begin{supertabular}{c@{$\;$}|p{.15\linewidth}@{}p{.15\linewidth}p{.15\linewidth}p{.15\linewidth}p{.15\linewidth}p{.15\linewidth}}
    \# & $\;$A & \multicolumn{4}{c}{Game Master} & $\;\:$B\\
    \hline

    \theutterance \stepcounter{utterance}  
    & & & \multicolumn{2}{p{0.3\linewidth}}{
        \cellcolor[rgb]{0.9,0.9,0.9}{
            \makecell[{{p{\linewidth}}}]{
                \texttt{\tiny{[GM$|$GM]}}
                {'hidden': {'team': ['tīmeklis', 'šūna', 'gaiss', 'ūdens', 'telpa', 'princese', 'krāns', 'apmetnis', 'adata'], 'innocent': ['ķengurs', 'dzelzs', 'straume', 'lineāls', 'ceļojums', 'mērogs', 'inde'], 'opponent': ['līnija', 'ronis', 'apakšdaļa', 'amerika', 'pele', 'bedre', 'skrūve', 'ķīna'], 'assassin': ['krēsls']}, 'revealed': {'team': {'team': [], 'innocent': [], 'opponent': [], 'assassin': []}, 'opponent': {'team': [], 'innocent': [], 'opponent': [], 'assassin': []}}}
            }
        }
    }
    & & \\ \\

    \theutterance \stepcounter{utterance}  
    & & \multicolumn{4}{p{0.6\linewidth}}{
        \cellcolor[rgb]{0.9,0.9,0.9}{
            \makecell[{{p{\linewidth}}}]{
                \texttt{\tiny{[P1$\langle$GM]}}
                \texttt{Sniedziet vienu vienvārdu norādi, kas ir saistīta ar kādu no jūsu komandas vārdiem no tālāk dotā saraksta. Jūs varat izvēlēties, uz cik un kādiem vārdiem attiecas jūsu norāde.} \\
\texttt{Mērķis ir panākt, lai jūsu partneris uzmin visus jūsu komandas vārdus, pamatojoties uz jūsu dotajām norādēm, pirms otra komanda ir atklājusi visus savus vārdus. Ja jūsu partneris uzmin "pretinieka" vai "novirzītāja" vārdu, viņa uzminēšanas kārta beidzas, un pretinieku komanda var veikt savu kārtu. Ja jūsu partneris uzmin "nogalināšanas" vārdu, spēle beidzas un jūsu komanda uzreiz zaudē.} \\
\texttt{Jo vairāk vārdu, uz kuriem attiecas jūsu norāde, jo ātrāk jūs varat uzvarēt spēli, bet jo grūtāk asociācijas var būt uzminēt jūsu partnerim, tāpēc izvēlieties savu norādes vārdu gudri. Norādes vārdam ir jābūt semantiski saistītam ar mērķa vārdiem, tas nedrīkst būt viens no sarakstos iekļautajiem vārdiem vai saturēt to daļas.} \\
\\ 
\texttt{Vienvārdu norādi un ar komatu atdalīto saistīto mērķa vārdu sarakstu vienmēr sniedziet šādā formātā un atbildi veidojiet pēc iespējas īsāku, nekad neiekļaujiet citu tekstu, kas nav prasīts šajā formā:} \\
\\ 
\texttt{PADOMS: <VĀRDS>} \\
\texttt{MĒRĶI: <MĒRĶI>} \\
\\ 
\texttt{Jūsu komandas vārdi ir: tīmeklis, šūna, gaiss, ūdens, telpa, princese, krāns, apmetnis, adata.} \\
\texttt{Jūsu pretinieka komandas vārdi ir: līnija, ronis, apakšdaļa, amerika, pele, bedre, skrūve, ķīna.} \\
\texttt{Novirzītāja vārdi ir: ķengurs, dzelzs, straume, lineāls, ceļojums, mērogs, inde.} \\
\texttt{Nogalināšanas vārdi ir: krēsls.} \\
            }
        }
    }
    & \\ \\

    \theutterance \stepcounter{utterance}  
    & \multicolumn{4}{p{0.6\linewidth}}{
        \cellcolor[rgb]{0.9,0.9,0.9}{
            \makecell[{{p{\linewidth}}}]{
                \texttt{\tiny{[P1$\rangle$GM]}}
                \texttt{PADOMS: dabi} \\
\texttt{MĒRĶI: gaiss, ūdens, telpa} \\
            }
        }
    }
    & & \\ \\

    \theutterance \stepcounter{utterance}  
    & & & \multicolumn{2}{p{0.3\linewidth}}{
        \cellcolor[rgb]{0.9,0.9,0.9}{
            \makecell[{{p{\linewidth}}}]{
                \texttt{\tiny{[GM$|$GM]}}
                \texttt{dabi} \\
            }
        }
    }
    & & \\ \\

    \theutterance \stepcounter{utterance}  
    & & & \multicolumn{2}{p{0.3\linewidth}}{
        \cellcolor[rgb]{0.9,0.9,0.9}{
            \makecell[{{p{\linewidth}}}]{
                \texttt{\tiny{[GM$|$GM]}}
                ['gaiss', 'ūdens', 'telpa']
            }
        }
    }
    & & \\ \\

    \theutterance \stepcounter{utterance}  
    & & & \multicolumn{2}{p{0.3\linewidth}}{
        \cellcolor[rgb]{0.9,0.9,0.9}{
            \makecell[{{p{\linewidth}}}]{
                \texttt{\tiny{[GM$|$GM]}}
                {'word': 'gaiss', 'assignment': 'team'}
            }
        }
    }
    & & \\ \\

    \theutterance \stepcounter{utterance}  
    & & & \multicolumn{2}{p{0.3\linewidth}}{
        \cellcolor[rgb]{0.9,0.9,0.9}{
            \makecell[{{p{\linewidth}}}]{
                \texttt{\tiny{[GM$|$GM]}}
                {'word': 'ūdens', 'assignment': 'team'}
            }
        }
    }
    & & \\ \\

    \theutterance \stepcounter{utterance}  
    & & & \multicolumn{2}{p{0.3\linewidth}}{
        \cellcolor[rgb]{0.9,0.9,0.9}{
            \makecell[{{p{\linewidth}}}]{
                \texttt{\tiny{[GM$|$GM]}}
                {'word': 'telpa', 'assignment': 'team'}
            }
        }
    }
    & & \\ \\

    \theutterance \stepcounter{utterance}  
    & & \multicolumn{4}{p{0.6\linewidth}}{
        \cellcolor[rgb]{0.9,0.9,0.9}{
            \makecell[{{p{\linewidth}}}]{
                \texttt{\tiny{[GM$\rangle$P2]}}
                \texttt{Sniedziet ar komatiem atdalītu sarakstu ar līdz 3 vārdiem no turpmāk minētā saraksta, kas vislabāk attiecas uz vārdu 'dabi' vai ir visciešāk saistīti ar to. Vienmēr sāciet minējumu sarakstu ar MINĒJUMS: un atbildē neiekļaujiet nekādu citu tekstu.} \\
\\ 
\texttt{mērogs, apmetnis, ūdens, līnija, ķīna, gaiss, pele, ķengurs, telpa, adata, skrūve, apakšdaļa, tīmeklis, ceļojums, krāns, šūna, dzelzs, straume, bedre, amerika, princese, lineāls, krēsls, inde, ronis} \\
            }
        }
    }
    & \\ \\

    \theutterance \stepcounter{utterance}  
    & & & \multicolumn{4}{p{0.6\linewidth}}{
        \cellcolor[rgb]{0.9,0.9,0.9}{
            \makecell[{{p{\linewidth}}}]{
                \texttt{\tiny{[GM$\langle$P2]}}
                \texttt{MINĒJUMS: telpa, tīmeklis, līnija} \\
            }
        }
    }
     \\ \\

    \theutterance \stepcounter{utterance}  
    & & & \multicolumn{2}{p{0.3\linewidth}}{
        \cellcolor[rgb]{0.9,0.9,0.9}{
            \makecell[{{p{\linewidth}}}]{
                \texttt{\tiny{[GM$|$GM]}}
                ['telpa', 'tīmeklis', 'līnija']
            }
        }
    }
    & & \\ \\

    \theutterance \stepcounter{utterance}  
    & & & \multicolumn{2}{p{0.3\linewidth}}{
        \cellcolor[rgb]{0.9,0.9,0.9}{
            \makecell[{{p{\linewidth}}}]{
                \texttt{\tiny{[GM$|$GM]}}
                {'word': 'telpa', 'assignment': 'team'}
            }
        }
    }
    & & \\ \\

    \theutterance \stepcounter{utterance}  
    & & & \multicolumn{2}{p{0.3\linewidth}}{
        \cellcolor[rgb]{0.9,0.9,0.9}{
            \makecell[{{p{\linewidth}}}]{
                \texttt{\tiny{[GM$|$GM]}}
                {'word': 'telpa', 'assignment': 'team'}
            }
        }
    }
    & & \\ \\

    \theutterance \stepcounter{utterance}  
    & & & \multicolumn{2}{p{0.3\linewidth}}{
        \cellcolor[rgb]{0.9,0.9,0.9}{
            \makecell[{{p{\linewidth}}}]{
                \texttt{\tiny{[GM$|$GM]}}
                {'word': 'tīmeklis', 'assignment': 'team'}
            }
        }
    }
    & & \\ \\

    \theutterance \stepcounter{utterance}  
    & & & \multicolumn{2}{p{0.3\linewidth}}{
        \cellcolor[rgb]{0.9,0.9,0.9}{
            \makecell[{{p{\linewidth}}}]{
                \texttt{\tiny{[GM$|$GM]}}
                {'word': 'līnija', 'assignment': 'opponent'}
            }
        }
    }
    & & \\ \\

    \theutterance \stepcounter{utterance}  
    & & & \multicolumn{2}{p{0.3\linewidth}}{
        \cellcolor[rgb]{0.9,0.9,0.9}{
            \makecell[{{p{\linewidth}}}]{
                \texttt{\tiny{[GM$|$GM]}}
                \texttt{līnija} \\
            }
        }
    }
    & & \\ \\

    \theutterance \stepcounter{utterance}  
    & & & \multicolumn{2}{p{0.3\linewidth}}{
        \cellcolor[rgb]{0.9,0.9,0.9}{
            \makecell[{{p{\linewidth}}}]{
                \texttt{\tiny{[GM$|$GM]}}
                {'word': 'ķīna', 'assignment': 'opponent'}
            }
        }
    }
    & & \\ \\

    \theutterance \stepcounter{utterance}  
    & & & \multicolumn{2}{p{0.3\linewidth}}{
        \cellcolor[rgb]{0.9,0.9,0.9}{
            \makecell[{{p{\linewidth}}}]{
                \texttt{\tiny{[GM$|$GM]}}
                {'word': 'apakšdaļa', 'assignment': 'opponent'}
            }
        }
    }
    & & \\ \\

    \theutterance \stepcounter{utterance}  
    & & & \multicolumn{2}{p{0.3\linewidth}}{
        \cellcolor[rgb]{0.9,0.9,0.9}{
            \makecell[{{p{\linewidth}}}]{
                \texttt{\tiny{[GM$|$GM]}}
                {'hidden': {'team': ['šūna', 'gaiss', 'ūdens', 'princese', 'krāns', 'apmetnis', 'adata'], 'innocent': ['ķengurs', 'dzelzs', 'straume', 'lineāls', 'ceļojums', 'mērogs', 'inde'], 'opponent': ['ronis', 'amerika', 'pele', 'bedre', 'skrūve'], 'assassin': ['krēsls']}, 'revealed': {'team': {'team': ['telpa', 'tīmeklis'], 'innocent': [], 'opponent': ['līnija'], 'assassin': []}, 'opponent': {'team': [], 'innocent': [], 'opponent': ['ķīna', 'apakšdaļa'], 'assassin': []}}}
            }
        }
    }
    & & \\ \\

    \theutterance \stepcounter{utterance}  
    & & \multicolumn{4}{p{0.6\linewidth}}{
        \cellcolor[rgb]{0.9,0.9,0.9}{
            \makecell[{{p{\linewidth}}}]{
                \texttt{\tiny{[P1$\langle$GM]}}
                \texttt{Vārdi telpa, tīmeklis tika uzminēti pareizi. Vārds līnija tika uzminēts, bet tas ir pretinieks vārds. Jūsu komandas biedra gājiens beidzās tur.} \\
\texttt{Tagad sniedziet vēl vienu norādi, kas attiecas uz dažiem no jūsu atlikušajiem komandas vārdiem, un saistīto mērķa vārdu sarakstu. Neaizmirstiet sākt savu norādi ar 'PADOMS: ', pēc tam ievietot jaunu rindu un sākt ar komatiem atdalīto mērķa vārdu sarakstu ar 'MĒRĶI: '. Piezīme: salīdzinājumā ar iepriekšējiem pieprasījumiem daži vārdi no sarakstiem ir izņemti.} \\
\\ 
\texttt{Jūsu atlikušie komandas vārdi ir: šūna, gaiss, ūdens, princese, krāns, apmetnis, adata.} \\
\texttt{Atlikušie pretinieka vārdi ir: ronis, apakšdaļa, amerika, pele, bedre, skrūve, ķīna.} \\
\texttt{Atlikušie traucētājvārdi ir: ķengurs, dzelzs, straume, lineāls, ceļojums, mērogs, inde.} \\
\texttt{Atlikušie slepkavības vārdi ir: krēsls.} \\
            }
        }
    }
    & \\ \\

    \theutterance \stepcounter{utterance}  
    & \multicolumn{4}{p{0.6\linewidth}}{
        \cellcolor[rgb]{0.9,0.9,0.9}{
            \makecell[{{p{\linewidth}}}]{
                \texttt{\tiny{[P1$\rangle$GM]}}
                \texttt{PADOMS: šūšanas} \\
\texttt{MĒRĶI: šūna, adata, princese} \\
            }
        }
    }
    & & \\ \\

    \theutterance \stepcounter{utterance}  
    & & & \multicolumn{2}{p{0.3\linewidth}}{
        \cellcolor[rgb]{0.9,0.9,0.9}{
            \makecell[{{p{\linewidth}}}]{
                \texttt{\tiny{[GM$|$GM]}}
                \texttt{šūšanas} \\
            }
        }
    }
    & & \\ \\

    \theutterance \stepcounter{utterance}  
    & & & \multicolumn{2}{p{0.3\linewidth}}{
        \cellcolor[rgb]{0.9,0.9,0.9}{
            \makecell[{{p{\linewidth}}}]{
                \texttt{\tiny{[GM$|$GM]}}
                ['šūna', 'adata', 'princese']
            }
        }
    }
    & & \\ \\

    \theutterance \stepcounter{utterance}  
    & & & \multicolumn{2}{p{0.3\linewidth}}{
        \cellcolor[rgb]{0.9,0.9,0.9}{
            \makecell[{{p{\linewidth}}}]{
                \texttt{\tiny{[GM$|$GM]}}
                {'word': 'šūna', 'assignment': 'team'}
            }
        }
    }
    & & \\ \\

    \theutterance \stepcounter{utterance}  
    & & & \multicolumn{2}{p{0.3\linewidth}}{
        \cellcolor[rgb]{0.9,0.9,0.9}{
            \makecell[{{p{\linewidth}}}]{
                \texttt{\tiny{[GM$|$GM]}}
                {'word': 'adata', 'assignment': 'team'}
            }
        }
    }
    & & \\ \\

    \theutterance \stepcounter{utterance}  
    & & & \multicolumn{2}{p{0.3\linewidth}}{
        \cellcolor[rgb]{0.9,0.9,0.9}{
            \makecell[{{p{\linewidth}}}]{
                \texttt{\tiny{[GM$|$GM]}}
                {'word': 'princese', 'assignment': 'team'}
            }
        }
    }
    & & \\ \\

    \theutterance \stepcounter{utterance}  
    & & \multicolumn{4}{p{0.6\linewidth}}{
        \cellcolor[rgb]{0.9,0.9,0.9}{
            \makecell[{{p{\linewidth}}}]{
                \texttt{\tiny{[GM$\rangle$P2]}}
                \texttt{Vārdi telpa, tīmeklis tika uzminēti pareizi. Vārds līnija tika uzminēts, bet tas ir pretinieks vārds. Jūsu gājiens beidzās tur.} \\
\texttt{Tagad sniedziet vēl vienu ar komatiem atdalītu sarakstu, kurā ir vismaz 1 un ne vairāk kā 3 vārdi no turpmāk minētā vārdu saraksta, kas vislabāk sasaucas vai ir visciešāk saistīti ar vārdu 'šūšanas'. Neaizmirstiet atbildi sākt ar MINĒJUMS:. Piezīme: salīdzinājumā ar iepriekšējiem pieprasījumiem daži vārdi no saraksta ir izņemti.} \\
\\ 
\texttt{mērogs, apmetnis, ūdens, gaiss, pele, ķengurs, adata, skrūve, ceļojums, krāns, šūna, dzelzs, straume, bedre, amerika, princese, lineāls, krēsls, inde, ronis} \\
            }
        }
    }
    & \\ \\

    \theutterance \stepcounter{utterance}  
    & & & \multicolumn{4}{p{0.6\linewidth}}{
        \cellcolor[rgb]{0.9,0.9,0.9}{
            \makecell[{{p{\linewidth}}}]{
                \texttt{\tiny{[GM$\langle$P2]}}
                \texttt{MINĒJUMS: adata, bedre, lineāls} \\
            }
        }
    }
     \\ \\

    \theutterance \stepcounter{utterance}  
    & & & \multicolumn{2}{p{0.3\linewidth}}{
        \cellcolor[rgb]{0.9,0.9,0.9}{
            \makecell[{{p{\linewidth}}}]{
                \texttt{\tiny{[GM$|$GM]}}
                ['adata', 'bedre', 'lineāls']
            }
        }
    }
    & & \\ \\

    \theutterance \stepcounter{utterance}  
    & & & \multicolumn{2}{p{0.3\linewidth}}{
        \cellcolor[rgb]{0.9,0.9,0.9}{
            \makecell[{{p{\linewidth}}}]{
                \texttt{\tiny{[GM$|$GM]}}
                {'word': 'adata', 'assignment': 'team'}
            }
        }
    }
    & & \\ \\

    \theutterance \stepcounter{utterance}  
    & & & \multicolumn{2}{p{0.3\linewidth}}{
        \cellcolor[rgb]{0.9,0.9,0.9}{
            \makecell[{{p{\linewidth}}}]{
                \texttt{\tiny{[GM$|$GM]}}
                {'word': 'adata', 'assignment': 'team'}
            }
        }
    }
    & & \\ \\

    \theutterance \stepcounter{utterance}  
    & & & \multicolumn{2}{p{0.3\linewidth}}{
        \cellcolor[rgb]{0.9,0.9,0.9}{
            \makecell[{{p{\linewidth}}}]{
                \texttt{\tiny{[GM$|$GM]}}
                {'word': 'bedre', 'assignment': 'opponent'}
            }
        }
    }
    & & \\ \\

    \theutterance \stepcounter{utterance}  
    & & & \multicolumn{2}{p{0.3\linewidth}}{
        \cellcolor[rgb]{0.9,0.9,0.9}{
            \makecell[{{p{\linewidth}}}]{
                \texttt{\tiny{[GM$|$GM]}}
                \texttt{bedre} \\
            }
        }
    }
    & & \\ \\

    \theutterance \stepcounter{utterance}  
    & & & \multicolumn{2}{p{0.3\linewidth}}{
        \cellcolor[rgb]{0.9,0.9,0.9}{
            \makecell[{{p{\linewidth}}}]{
                \texttt{\tiny{[GM$|$GM]}}
                {'word': 'amerika', 'assignment': 'opponent'}
            }
        }
    }
    & & \\ \\

    \theutterance \stepcounter{utterance}  
    & & & \multicolumn{2}{p{0.3\linewidth}}{
        \cellcolor[rgb]{0.9,0.9,0.9}{
            \makecell[{{p{\linewidth}}}]{
                \texttt{\tiny{[GM$|$GM]}}
                {'word': 'pele', 'assignment': 'opponent'}
            }
        }
    }
    & & \\ \\

    \theutterance \stepcounter{utterance}  
    & & & \multicolumn{2}{p{0.3\linewidth}}{
        \cellcolor[rgb]{0.9,0.9,0.9}{
            \makecell[{{p{\linewidth}}}]{
                \texttt{\tiny{[GM$|$GM]}}
                {'hidden': {'team': ['šūna', 'gaiss', 'ūdens', 'princese', 'krāns', 'apmetnis'], 'innocent': ['ķengurs', 'dzelzs', 'straume', 'lineāls', 'ceļojums', 'mērogs', 'inde'], 'opponent': ['ronis', 'skrūve'], 'assassin': ['krēsls']}, 'revealed': {'team': {'team': ['telpa', 'tīmeklis', 'adata'], 'innocent': [], 'opponent': ['līnija', 'bedre'], 'assassin': []}, 'opponent': {'team': [], 'innocent': [], 'opponent': ['ķīna', 'apakšdaļa', 'amerika', 'pele'], 'assassin': []}}}
            }
        }
    }
    & & \\ \\

    \theutterance \stepcounter{utterance}  
    & & \multicolumn{4}{p{0.6\linewidth}}{
        \cellcolor[rgb]{0.9,0.9,0.9}{
            \makecell[{{p{\linewidth}}}]{
                \texttt{\tiny{[P1$\langle$GM]}}
                \texttt{Vārds adata tika uzminēts pareizi. Vārds bedre tika uzminēts, bet tas ir pretinieks vārds. Jūsu komandas biedra gājiens beidzās tur.} \\
\texttt{Tagad sniedziet vēl vienu norādi, kas attiecas uz dažiem no jūsu atlikušajiem komandas vārdiem, un saistīto mērķa vārdu sarakstu. Neaizmirstiet sākt savu norādi ar 'PADOMS: ', pēc tam ievietot jaunu rindu un sākt ar komatiem atdalīto mērķa vārdu sarakstu ar 'MĒRĶI: '. Piezīme: salīdzinājumā ar iepriekšējiem pieprasījumiem daži vārdi no sarakstiem ir izņemti.} \\
\\ 
\texttt{Jūsu atlikušie komandas vārdi ir: šūna, gaiss, ūdens, princese, krāns, apmetnis.} \\
\texttt{Atlikušie pretinieka vārdi ir: ronis, amerika, pele, skrūve.} \\
\texttt{Atlikušie traucētājvārdi ir: ķengurs, dzelzs, straume, lineāls, ceļojums, mērogs, inde.} \\
\texttt{Atlikušie slepkavības vārdi ir: krēsls.} \\
            }
        }
    }
    & \\ \\

    \theutterance \stepcounter{utterance}  
    & \multicolumn{4}{p{0.6\linewidth}}{
        \cellcolor[rgb]{0.9,0.9,0.9}{
            \makecell[{{p{\linewidth}}}]{
                \texttt{\tiny{[P1$\rangle$GM]}}
                \texttt{PADOMS: plūsma} \\
\texttt{MĒRĶI: gaiss, ūdens, krāns} \\
            }
        }
    }
    & & \\ \\

    \theutterance \stepcounter{utterance}  
    & & & \multicolumn{2}{p{0.3\linewidth}}{
        \cellcolor[rgb]{0.9,0.9,0.9}{
            \makecell[{{p{\linewidth}}}]{
                \texttt{\tiny{[GM$|$GM]}}
                \texttt{plūsma} \\
            }
        }
    }
    & & \\ \\

    \theutterance \stepcounter{utterance}  
    & & & \multicolumn{2}{p{0.3\linewidth}}{
        \cellcolor[rgb]{0.9,0.9,0.9}{
            \makecell[{{p{\linewidth}}}]{
                \texttt{\tiny{[GM$|$GM]}}
                ['gaiss', 'ūdens', 'krāns']
            }
        }
    }
    & & \\ \\

    \theutterance \stepcounter{utterance}  
    & & & \multicolumn{2}{p{0.3\linewidth}}{
        \cellcolor[rgb]{0.9,0.9,0.9}{
            \makecell[{{p{\linewidth}}}]{
                \texttt{\tiny{[GM$|$GM]}}
                {'word': 'gaiss', 'assignment': 'team'}
            }
        }
    }
    & & \\ \\

    \theutterance \stepcounter{utterance}  
    & & & \multicolumn{2}{p{0.3\linewidth}}{
        \cellcolor[rgb]{0.9,0.9,0.9}{
            \makecell[{{p{\linewidth}}}]{
                \texttt{\tiny{[GM$|$GM]}}
                {'word': 'ūdens', 'assignment': 'team'}
            }
        }
    }
    & & \\ \\

    \theutterance \stepcounter{utterance}  
    & & & \multicolumn{2}{p{0.3\linewidth}}{
        \cellcolor[rgb]{0.9,0.9,0.9}{
            \makecell[{{p{\linewidth}}}]{
                \texttt{\tiny{[GM$|$GM]}}
                {'word': 'krāns', 'assignment': 'team'}
            }
        }
    }
    & & \\ \\

    \theutterance \stepcounter{utterance}  
    & & \multicolumn{4}{p{0.6\linewidth}}{
        \cellcolor[rgb]{0.9,0.9,0.9}{
            \makecell[{{p{\linewidth}}}]{
                \texttt{\tiny{[GM$\rangle$P2]}}
                \texttt{Vārds adata tika uzminēts pareizi. Vārds bedre tika uzminēts, bet tas ir pretinieks vārds. Jūsu gājiens beidzās tur.} \\
\texttt{Tagad sniedziet vēl vienu ar komatiem atdalītu sarakstu, kurā ir vismaz 1 un ne vairāk kā 3 vārdi no turpmāk minētā vārdu saraksta, kas vislabāk sasaucas vai ir visciešāk saistīti ar vārdu 'plūsma'. Neaizmirstiet atbildi sākt ar MINĒJUMS:. Piezīme: salīdzinājumā ar iepriekšējiem pieprasījumiem daži vārdi no saraksta ir izņemti.} \\
\\ 
\texttt{mērogs, apmetnis, ūdens, gaiss, ķengurs, skrūve, ceļojums, krāns, šūna, dzelzs, straume, princese, lineāls, krēsls, inde, ronis} \\
            }
        }
    }
    & \\ \\

    \theutterance \stepcounter{utterance}  
    & & & \multicolumn{4}{p{0.6\linewidth}}{
        \cellcolor[rgb]{0.9,0.9,0.9}{
            \makecell[{{p{\linewidth}}}]{
                \texttt{\tiny{[GM$\langle$P2]}}
                \texttt{MINĒJUMS: ūdens, straume, gaiss} \\
            }
        }
    }
     \\ \\

    \theutterance \stepcounter{utterance}  
    & & & \multicolumn{2}{p{0.3\linewidth}}{
        \cellcolor[rgb]{0.9,0.9,0.9}{
            \makecell[{{p{\linewidth}}}]{
                \texttt{\tiny{[GM$|$GM]}}
                ['ūdens', 'straume', 'gaiss']
            }
        }
    }
    & & \\ \\

    \theutterance \stepcounter{utterance}  
    & & & \multicolumn{2}{p{0.3\linewidth}}{
        \cellcolor[rgb]{0.9,0.9,0.9}{
            \makecell[{{p{\linewidth}}}]{
                \texttt{\tiny{[GM$|$GM]}}
                {'word': 'ūdens', 'assignment': 'team'}
            }
        }
    }
    & & \\ \\

    \theutterance \stepcounter{utterance}  
    & & & \multicolumn{2}{p{0.3\linewidth}}{
        \cellcolor[rgb]{0.9,0.9,0.9}{
            \makecell[{{p{\linewidth}}}]{
                \texttt{\tiny{[GM$|$GM]}}
                {'word': 'ūdens', 'assignment': 'team'}
            }
        }
    }
    & & \\ \\

    \theutterance \stepcounter{utterance}  
    & & & \multicolumn{2}{p{0.3\linewidth}}{
        \cellcolor[rgb]{0.9,0.9,0.9}{
            \makecell[{{p{\linewidth}}}]{
                \texttt{\tiny{[GM$|$GM]}}
                {'word': 'straume', 'assignment': 'innocent'}
            }
        }
    }
    & & \\ \\

    \theutterance \stepcounter{utterance}  
    & & & \multicolumn{2}{p{0.3\linewidth}}{
        \cellcolor[rgb]{0.9,0.9,0.9}{
            \makecell[{{p{\linewidth}}}]{
                \texttt{\tiny{[GM$|$GM]}}
                \texttt{straume} \\
            }
        }
    }
    & & \\ \\

    \theutterance \stepcounter{utterance}  
    & & & \multicolumn{2}{p{0.3\linewidth}}{
        \cellcolor[rgb]{0.9,0.9,0.9}{
            \makecell[{{p{\linewidth}}}]{
                \texttt{\tiny{[GM$|$GM]}}
                {'word': 'ronis', 'assignment': 'opponent'}
            }
        }
    }
    & & \\ \\

    \theutterance \stepcounter{utterance}  
    & & & \multicolumn{2}{p{0.3\linewidth}}{
        \cellcolor[rgb]{0.9,0.9,0.9}{
            \makecell[{{p{\linewidth}}}]{
                \texttt{\tiny{[GM$|$GM]}}
                {'word': 'skrūve', 'assignment': 'opponent'}
            }
        }
    }
    & & \\ \\

    \theutterance \stepcounter{utterance}  
    & & & \multicolumn{2}{p{0.3\linewidth}}{
        \cellcolor[rgb]{0.9,0.9,0.9}{
            \makecell[{{p{\linewidth}}}]{
                \texttt{\tiny{[GM$|$GM]}}
                {'hidden': {'team': ['šūna', 'gaiss', 'princese', 'krāns', 'apmetnis'], 'innocent': ['ķengurs', 'dzelzs', 'lineāls', 'ceļojums', 'mērogs', 'inde'], 'opponent': [], 'assassin': ['krēsls']}, 'revealed': {'team': {'team': ['telpa', 'tīmeklis', 'adata', 'ūdens'], 'innocent': ['straume'], 'opponent': ['līnija', 'bedre'], 'assassin': []}, 'opponent': {'team': [], 'innocent': [], 'opponent': ['ķīna', 'apakšdaļa', 'amerika', 'pele', 'ronis', 'skrūve'], 'assassin': []}}}
            }
        }
    }
    & & \\ \\

    \theutterance \stepcounter{utterance}  
    & & \multicolumn{4}{p{0.6\linewidth}}{
        \cellcolor[rgb]{0.9,0.9,0.9}{
            \makecell[{{p{\linewidth}}}]{
                \texttt{\tiny{[P1$\langle$GM]}}
                \texttt{Vārds ūdens tika uzminēts pareizi. Vārds straume tika uzminēts, bet tas ir nevainojams vārds. Jūsu komandas biedra gājiens beidzās tur.} \\
\texttt{Tagad sniedziet vēl vienu norādi, kas attiecas uz dažiem no jūsu atlikušajiem komandas vārdiem, un saistīto mērķa vārdu sarakstu. Neaizmirstiet sākt savu norādi ar 'PADOMS: ', pēc tam ievietot jaunu rindu un sākt ar komatiem atdalīto mērķa vārdu sarakstu ar 'MĒRĶI: '. Piezīme: salīdzinājumā ar iepriekšējiem pieprasījumiem daži vārdi no sarakstiem ir izņemti.} \\
\\ 
\texttt{Jūsu atlikušie komandas vārdi ir: šūna, gaiss, princese, krāns, apmetnis.} \\
\texttt{Atlikušie pretinieka vārdi ir: ronis, skrūve.} \\
\texttt{Atlikušie traucētājvārdi ir: ķengurs, dzelzs, lineāls, ceļojums, mērogs, inde.} \\
\texttt{Atlikušie slepkavības vārdi ir: krēsls.} \\
            }
        }
    }
    & \\ \\

    \theutterance \stepcounter{utterance}  
    & \multicolumn{4}{p{0.6\linewidth}}{
        \cellcolor[rgb]{0.9,0.9,0.9}{
            \makecell[{{p{\linewidth}}}]{
                \texttt{\tiny{[P1$\rangle$GM]}}
                \texttt{PADOMS: sēde} \\
\texttt{MĒRĶI: šūna, apmetnis} \\
            }
        }
    }
    & & \\ \\

    \theutterance \stepcounter{utterance}  
    & & & \multicolumn{2}{p{0.3\linewidth}}{
        \cellcolor[rgb]{0.9,0.9,0.9}{
            \makecell[{{p{\linewidth}}}]{
                \texttt{\tiny{[GM$|$GM]}}
                \texttt{sēde} \\
            }
        }
    }
    & & \\ \\

    \theutterance \stepcounter{utterance}  
    & & & \multicolumn{2}{p{0.3\linewidth}}{
        \cellcolor[rgb]{0.9,0.9,0.9}{
            \makecell[{{p{\linewidth}}}]{
                \texttt{\tiny{[GM$|$GM]}}
                ['šūna', 'apmetnis']
            }
        }
    }
    & & \\ \\

    \theutterance \stepcounter{utterance}  
    & & & \multicolumn{2}{p{0.3\linewidth}}{
        \cellcolor[rgb]{0.9,0.9,0.9}{
            \makecell[{{p{\linewidth}}}]{
                \texttt{\tiny{[GM$|$GM]}}
                {'word': 'šūna', 'assignment': 'team'}
            }
        }
    }
    & & \\ \\

    \theutterance \stepcounter{utterance}  
    & & & \multicolumn{2}{p{0.3\linewidth}}{
        \cellcolor[rgb]{0.9,0.9,0.9}{
            \makecell[{{p{\linewidth}}}]{
                \texttt{\tiny{[GM$|$GM]}}
                {'word': 'apmetnis', 'assignment': 'team'}
            }
        }
    }
    & & \\ \\

    \theutterance \stepcounter{utterance}  
    & & & \multicolumn{2}{p{0.3\linewidth}}{
        \cellcolor[rgb]{0.9,0.9,0.9}{
            \makecell[{{p{\linewidth}}}]{
                \texttt{\tiny{[GM$|$GM]}}
                \texttt{opponent has won} \\
            }
        }
    }
    & & \\ \\

\end{supertabular}
}

\end{document}
